% RESEARCH_DRAFT_SHA256: 079ec3cf871c111321fddb5f3816044372ad0eefdf6c3e3d7fdce4ba7eb54aba
% Layout and source-reference wrappers expanded for the web renderer.
\documentclass[11pt]{article}
\usepackage[margin=25mm]{geometry}
\usepackage{amsthm}
\usepackage{fontspec,unicode-math}
\setmainfont{Latin Modern Roman}
\setmathfont{Latin Modern Math}
\usepackage{amsmath,amssymb,mathtools}
\usepackage{xurl,hyperref}
\hypersetup{colorlinks=true,urlcolor=blue,linkcolor=blue}
\setcounter{secnumdepth}{0}
\setlength{\parindent}{0pt}
\setlength{\parskip}{5pt}
\setlength{\emergencystretch}{3em}
\newcommand{\sref}[2]{\hyperlink{source:#1:#2}{[S#2]}}
\newcommand{\cataloguescope}[1]{\paragraph{Source question.}#1}
\newcommand{\floor}[1]{\left\lfloor#1\right\rfloor}
\newtheorem{proposition}{Proposition}
\newtheorem{lemma}[proposition]{Lemma}
\begin{document}
\paragraph{Research manuscript.} OpenAI \texttt{gpt-6-astra}, reasoning \texttt{ultra}; 11 October 2026.
\paragraph{Preserved source body.} The complete catalogue mathematical text and sources follow with layout and source links expanded, including source qualifications and its historical preparation notes. New research begins at the section headed Research attempt.

% UnsolvedMath ID 20001217; source catalogue code AIM-COMPUTATION-0055
\section{ML degree, feasibility, and a Gaussian likelihood-root trichotomy}\label{20001217}
{\small UnsolvedMath ID 20001217 · Computer Science}
\subsection{Definitions and mathematical statement}

\paragraph{Shared notation.}Standard number systems are denoted by $\mathbb N,\mathbb Z,\mathbb Q,\mathbb R,\mathbb C$; index ranges, computational models and specialized notation follow the cited problem source and the definitions below.


For a statistical model with parameter $\theta$ and observations $D$, the likelihood is $L(\theta;D)=\Pr_\theta(D)$, or the corresponding density. Maximum likelihood seeks a global maximum over the admissible parameter space. An algebraic likelihood equation system can have complex solutions or real solutions outside the statistical parameter domain. Valid covariance matrices must be positive definite, and valid cell probabilities must be nonnegative and normalized. Counting all algebraic roots is distinct from counting admissible likelihood modes.
\paragraph{Statement recorded in the supplied source.}
Question 41 (Ho¸ sten). How many complex roots to the likelihood equations can there be? What is the statistical significance of complex roots/ negative real roots or other roots that do not yield a valid statistical model (e.g. covariance matrices that are not positive definite)? \href{https://aimath.org/WWN/compalgstat/compalgstat.pdf}{[S2]}
\paragraph{Source status.}
The supplied catalogue labels this record partially solved. Its source-status metadata was checked against the supplied snapshot on 2026-10-11. This import does not independently establish current open status or mathematical correctness.
\subsection{Short English statement}
How many algebraic likelihood roots are possible, and what can inadmissible roots tell us statistically?
\subsection{Sources}
\begin{itemize}
\item[S1] Ulam AI and the UnsolvedMath Contributors, UnsolvedMath: A Curated Collection of Open Mathematics Problems, record 20001217 (AIM-COMPUTATION-0055). Catalogue curation and metadata: CC BY 4.0. The supplied source question is retained; definitions and the English summary are editorial additions.
\url{https://huggingface.co/datasets/ulamai/UnsolvedMath}.
\item[S2] AIM Workshop Problem Lists, original problem source linked by catalogue record AIM-COMPUTATION-0055
\url{https://aimath.org/WWN/compalgstat/compalgstat.pdf}.
\end{itemize}
\paragraph{Preparation and review.}
Original question attributed to the cited source and UnsolvedMath contributors. Codex (GPT-6; exact serving revision not exposed) assisted with definitions, English summaries, source inspection and formatting. No new human mathematical authorship or independent proof verification is asserted. Canonical ID 20001217 and source code AIM-COMPUTATION-0055 checked against the complete supplied registry and existing repository statements; known flagged repetitions and obvious identical conjectures excluded. Full mathematical-equivalence/literature review has not been performed. Statement preparation and final review followed contribution prompts C and D; automated checks verify structure and identity, not open status.
\label{end:20001217}


\section{Research attempt}
\subsection{Literature and scope}
Question 41 concerns algebraic likelihood roots and statistical admissibility without a fixed model. I tested whether the sign of a real parameter root alone determines its statistical meaning. Two elementary Gaussian parameterizations show why the parameter-to-distribution map must be retained.
\begin{proposition}
For independent centered bivariate Gaussian observations with covariance $\Sigma(t)=\operatorname{diag}(t,t^2)$, $t>0$, and empirical second moments $s_1,s_2>0$, the cleared score equation has exactly two simple real roots, one admissible and one inadmissible. In contrast, a negative root for the parameterization of a scalar variance by $t^2$ can be fully admissible.
\end{proposition}
\begin{proof}
Ignoring constants and multiplying by the positive sample size does not change maximizers. The per-observation log likelihood in the first model is
\[
\ell(t)=-\tfrac12(3\log t+s_1/t+s_2/t^2),\qquad
\ell'(t)=-\frac{3t^2-s_1t-2s_2}{2t^3}.
\]
The roots are $(s_1\pm\sqrt{s_1^2+24s_2})/6$. They are distinct, nonzero, and have opposite signs. The positive root is the unique maximum: the derivative is positive before it and negative after it on $(0,\infty)$, and $\ell$ tends to $-\infty$ at both ends. At the negative root the first diagonal entry of $\Sigma$ is negative, so the covariance is not positive definite and no Gaussian law is represented. Since neither root is zero, clearing $t^3$ did not create either root.

For a scalar zero-mean Gaussian with variance $t^2$ and domain $t\in\mathbb R\setminus\{0\}$, second moment $s>0$ gives $\ell(t)=-\log|t|-s/(2t^2)$. The score is $-1/t+s/t^3$, whose roots $\pm\sqrt{s}$ both represent the same valid maximum-likelihood variance $s$. Thus negative coordinate values alone have no invariant statistical interpretation.
\end{proof}
\subsection{Verification and remaining gaps}
At $s_1=s_2=1$, the first score factors as $(t-1)(3t+2)$ and its roots are $1,-2/3$. The second example at $s=1$ has roots $\pm1$. These exact factorizations test the failed sign-only classification. They do not bound the complex roots for arbitrary models, analyze singular likelihood equations, or supply a general meaning for complex solutions. No novelty or general root-count theorem is claimed.

\section{References and source attribution}
The original question is credited to the AIM workshop list prepared by Seth Sullivant (2004), \url{https://aimath.org/WWN/compalgstat/compalgstat.pdf}. The source question and catalogue attribution are reproduced above. This list was accessed on 11 October 2026; it is historical evidence of the question, not a certification of present open status.


\paragraph{Integration review.} Codex (GPT-6; exact serving revision not exposed) reviewed the written arguments and expanded layout and source-reference wrappers. Repository metadata and the production web renderer were checked using Python 3.13 and JavaScript with MathJax 3.2.2. This model-assisted review is not independent human proof verification; no PDF compilation or formal proof checking is asserted.

\section{Completion Estimate}
\noindent\textbf{COMPLETION: 0\%}
No substantiated new progress toward the entire original target is claimed. The worked results and failed route are delimited above; a known result, illustrative model, or auxiliary obstruction is not counted as original-target completion.

\section{Final status and attribution}
\textbf{Attempt status: unresolved.} This describes this research attempt, not a new ruling on the source catalogue's status or on external literature claims.
The mathematics and this self-audit were generated by OpenAI \texttt{gpt-6-astra}, reasoning \texttt{ultra}, on 11 October 2026. Codex (GPT-6; exact serving revision not exposed) orchestrated the batch and contributed formatting/repository review. No human mathematical authorship, independent human proof checking, or formal verification is claimed. The historical preparation paragraphs in the reproduced catalogue body concern that body only.

Contribution prompt A was applied to the complete definition and supplied source record, followed by the adversarial checks of prompt B. Local reproduction files are \texttt{batch005/checks.py}, \texttt{computations.json}, and the per-ID audit JSON. They were not public links at preparation time. Numerical checks are finite and do not prove unrestricted results. No LaTeX engine was found locally; compilation was not performed by this research agent.

\end{document}
